\begin{helper}
Let \(W\) be a finite graph with activation--priority law \(\nu\) satisfying
the Bernoulli activation atom formula and the full priority lower-rectangle
law with parameter \(p\in[0,1]\).  Let \(C\) be a finite candidate index set,
with embedded candidate vertex \(e(x)\in W\) for each \(x\in C\).  Let
\(Q\subseteq W\) be the fixed local-coordinate set, and let
\(b_i(x)\) be the finitely many outside blocker coordinates attached to
candidate \(x\) whenever \(H_i(x)\) holds.  Assume every neighbor of every
embedded candidate is either in \(Q\) or is one of these listed blockers:
\[
  e(x)y\in E(W),\ x\in C
  \quad\Longrightarrow\quad
  y\in Q\ \hbox{or}\ \exists i\,(H_i(x)\hbox{ and }y=b_i(x)).
\]
Assume also \(e(x)\in Q\) for every \(x\in C\), and assume that each listed
blocker which is present for a candidate is an actual neighbor:
\[
  H_i(x),\ x\in C \quad\Longrightarrow\quad e(x)b_i(x)\in E(W),
\]
Then, for every support cell \(D\), the true strict endpoint survival event
for \(C\) has the same \(\nu\)-mass on \(D\) as the displayed
support-relative event in which local neighbors in \(Q\) are tested by the
strict inequalities \(\pi(y)<\pi(e(x))\), while each listed outside blocker
is excluded by the negated opposite test
\[
  \neg\bigl(b_i(x)\in A,\ \pi(e(x))<\pi(b_i(x)),\
  0\le \pi(b_i(x))\le1\bigr).
\]
That is, after intersection with \(D\),
\[
\{I_W(A,\pi)\cap e(C)\ne\emptyset\}
\]
and this displayed local/blocker event have equal measure.
\end{helper}
\begin{proof}
We first record the finite coordinate set to which
\(\noderef{SamplingSplitOutsideBlockerTargetEndpointEquivalence}\) will be
applied.  Put
\[
  Q^*=Q\cup\{b_i(x):x\in C,\ H_i(x)\}.
\]
If \(x\in C\) and \(y\) is adjacent to \(e(x)\), the cover hypothesis gives
either \(y\in Q\) or \(y=b_i(x)\) for some present blocker \(H_i(x)\).  Hence
\(y\in Q^*\).  The cited endpoint-equivalence node therefore identifies the
true endpoint event pointwise with the event that some \(x\in C\) is active
and beats every active adjacent coordinate in \(Q^*\).

We next compare this \(Q^*\)-local condition with the displayed condition for
a fixed candidate \(x\).  Coordinates in \(Q\) contribute exactly the local
tests appearing in the statement.  For a coordinate \(y\in Q^*\setminus Q\)
which is active and adjacent to \(e(x)\), it may have entered \(Q^*\) as a
listed blocker for some other candidate, but the cover hypothesis applied to
the present candidate \(x\) and the same adjacent coordinate \(y\) gives
\[
  y\in Q\quad\hbox{or}\quad y=b_i(x)\hbox{ for some }i\hbox{ with }H_i(x).
\]
The first alternative is already handled by the local tests, and in the
second alternative the new support condition says that this listed blocker is
indeed a genuine neighbor of \(e(x)\).  Thus the only nonlocal strict endpoint
tests that must be compared with the displayed event are precisely
\[
  b_i(x)\in A \quad\Longrightarrow\quad \pi(b_i(x))<\pi(e(x))
  \qquad(H_i(x)).
\]
It remains only to compare this strict implication with the displayed negated
opposite test.  This comparison is pointwise valid away from unit-support
failures and ties, and those exceptional sets have zero \(\nu\)-mass as follows.

The Bernoulli atom formula makes each activation atom finite, since
\(\operatorname{ofReal}\) of a real number is finite.  Thus
\noderef{SamplingPriorityPositiveUnitSupport}, applied with the given
lower-rectangle law, gives a conull set \(U\) on which
\(0<\pi(v)\le1\) for every \(v\in W\).  For each fixed activation
atom \(A_0\), the events
\[
 \{A=A_0,\ \pi(q)\in[0,1]\text{ for all }q\in A_0\}
 \quad\text{and}\quad \{A=A_0\}
\]
agree on \(U\), so they have equal measure.  This supplies the
active-coordinate cube-mass hypothesis of
\noderef{SamplingNonadjacentPairOrderedChamberDiagonalNull}.

For \(x\in C\) and \(H_i(x)\), the adjacency hypothesis and
irreflexivity give \(e(x)\ne b_i(x)\).  Apply that diagonal-null theorem
to these two coordinates, with the given Bernoulli and rectangle laws and
the cube-mass identity just established.  It says that
\[
 \{e(x),b_i(x)\in A,\ \pi(q)\in[0,1]\text{ for every }q\in A,
       \ \pi(e(x))=\pi(b_i(x))\}
\]
is null.  Take the union of these finitely many null sets and \(U^c\).
Outside this union, an active candidate and an active listed blocker have
distinct priorities.  No assertion about ties involving inactive candidates
is needed, since both survival events require an active candidate.
An active listed blocker also satisfies
\(0\le\pi(b_i(x))\le1\), and trichotomy of the two distinct real priorities
gives exactly one of
\[
  \pi(b_i(x))<\pi(e(x)),\qquad
  \pi(e(x))<\pi(b_i(x)).
\]
Consequently, for an active candidate, the displayed condition
\[
\neg\bigl(b_i(x)\in A,\ \pi(e(x))<\pi(b_i(x)),\
0\le\pi(b_i(x))\le1\bigr)
\]
is equivalent, off the null exceptional set, to the strict endpoint implication
\[
  b_i(x)\in A \quad\Longrightarrow\quad \pi(b_i(x))<\pi(e(x)).
\]

Combining the pointwise endpoint equivalence on \(Q^*\) with this finite
null-set conversion for the listed blockers shows that the symmetric difference
between the true endpoint survival event and the displayed local/blocker event
is contained in the union of the unit-support complement and finitely many
active--active candidate--blocker tie sets just described.
That union has \(\nu\)-measure \(0\).
Intersecting both events with the arbitrary cell \(D\) can only shrink the
symmetric difference, so the two intersections differ by a null set.  Their
\(\nu\)-masses are therefore equal.
\end{proof}
